\documentclass[a4paper]{IEEEtran}

\usepackage[
    backend=biber,
    style=ieee
]{biblatex}
\addbibresource{references.bib}

\usepackage{amsmath,amssymb,amsfonts, bm}
\usepackage{graphicx}
\usepackage[hidelinks]{hyperref}
\usepackage{cleveref}
\usepackage{booktabs}

\begin{document}

\title{Network Intrusion Detection}

\author{
    \IEEEauthorblockN{Sigurd Riseth}\\
    \IEEEauthorblockA{
        \textit{Department of Physics}\\
        \textit{University of Oslo}\\
        Oslo, Norway \\
        sigurri@uio.no
    }
}

\maketitle

\begin{abstract}
...
\end{abstract}

\begin{IEEEkeywords}
keyword1, keyword2, keyword3
\end{IEEEkeywords}

%\tableofcontents

\section{Introduction}

How effectively can different machine learning algorithms distinguish between
legitimate network traffic and malicious activity?

Dataset: CIC-IDS2017.

We compare the performance of logistic regression, random forest, XGBoost, and FFNN.

We investigate binary versus multi-class classification, class imbalance, precision-recall
for attack classes, false positive rates, feature importance, training time and computational
requirements.

Accuracy is potentially misleading in intrusion detection. A model that achieves
high accuracy but fails to detect rare attacks is not useful in practice.

As a benchmark, we use the 2021 study by \textcite{maseer_2021}, which compares ten different ML methods on the dataset.

\section{Methods}

\printbibliography

\appendix

\end{document}
